\documentclass[11pt]{article}
\title{Effective Superluminal Travel: Conditional Geometry, Source Exclusions, and a Gated Research Program}
\date{}
\author{SIGMA AI-assisted research synthesis}
\usepackage[T1]{fontenc}
\usepackage[utf8]{inputenc}
\usepackage[margin=1in]{geometry}
\usepackage{amsmath,amssymb,array,booktabs,longtable,pdflscape}
\usepackage{microtype}
\usepackage[hyphens]{url}
\usepackage{titlesec}
\emergencystretch=2em
\raggedbottom
\linespread{1.07}
\usepackage{xcolor}
\definecolor{blue}{HTML}{24485E}
\titleformat{\section}{\Large\bfseries\sffamily\color{blue}}{\thesection}{0.8em}{}
\titleformat{\subsection}{\large\bfseries\sffamily}{\thesubsection}{0.7em}{}
\setlength{\parskip}{.6em}
\setlength{\parindent}{0pt}
\begin{document}
\maketitle

\begin{abstract}
Can an effective spacetime shortcut remain coherent once its causal structure, source dynamics, and propagation limits are examined together? This theoretical study separates a prescribed geometric shortcut from a physically realizable transport mechanism. A smooth finite-duration metric permits an exterior-clock transit shorter than a matched Minkowski reference while its carrier remains locally timelike and a global temporal function excludes closed causal curves. Under Einstein gravity, however, the target requires negative normal energy and strict null-energy violation. The supplied analytical arguments exclude a canonical scalar and a regular, positive-energy single-scalar extension as sources.

A related eternal superluminal geometry also has incomplete forward axial null geodesics with divergent parallelly propagated curvature, preventing a regular extension through their affected endpoints. A preferred-time wider-cone free field offers a distinct conceptual comparison, but not an interacting transport theory. These results support a staged program of source, finite-duration, quantum, and control tests; they establish neither a working drive nor a universal prohibition on effective superluminal travel.
\end{abstract}
\section{Introduction and physical framework}
The answerable question is whether the specified geometric shortcut survives necessary tests of causal consistency, admissible sourcing, and regular propagation\textemdash{}not whether an arbitrarily chosen metric can be called a drive. Effective superluminality is defined by an exterior-clock transit shorter than the corresponding light transit in a matched, unmodified Minkowski geometry. It is distinct from a massive carrier leaving its local metric light cone.

The bounded contribution combines conditional geometric deductions with source-class exclusions and a restricted geodesic obstruction. A preferred-time wider-cone field supplies a conceptually different comparison, but its free-field status prevents an equal-depth assessment of interacting transport. The investigation therefore ends with explicit decision gates rather than a claim that either framework is physically viable.

\section{Models, methods, and evidence}
The saved analysis prescribed smooth unit-lapse, flat-spatial-slice geometries and examined their local tangent norms, exterior-clock transit protocols, and temporal functions. It treated the finite-duration pulse separately from the eternal constant-speed geometry. The geometric calculations used signature $(-,+,+,+)$; source calculations assumed Einstein gravity with zero cosmological constant and $G_{ab}=8\pi G T_{ab}/c^4$.

The source analysis used the Hamiltonian and momentum constraints, spatial evolution equations, and a separate frame-curvature calculation to examine the required density, stresses, fluxes, and null contractions. A spherical-wall estimate applied angular integration and Cauchy\textendash{}Schwarz. Initial-value uniqueness, canonical-scalar stress, and the quadratic Hamiltonian and principal symbol of a minimally coupled $P(X,\phi)$ scalar then supplied distinct formation and source-stability tests.

The propagation analysis solved axial affine equations using a translating Killing symmetry and distinguished a subluminal complete-ray integral from superluminal rays approaching characteristic separatrices. Parallelly propagated tidal curvature tested whether the latter endpoints were removable coordinate effects, while the backward family supplied a control. These were analytical arguments; the supplied evidence contains no completed numerical simulation, measured observation, or inspected domain-literature assessment.

The bounded evidence index omitted 7496 saved research fragments across 8 sources. Claims that could depend on those fragments remain provisional; the full omission ledger is retained with the archived research record.
\section{An operational shortcut without local superluminal motion}
A finite-duration geometry makes the proposed shortcut operationally precise: an exterior-clock transit can be shorter than a matched Minkowski light-transit time without the information carrier leaving its local photon cone. The construction separates three questions that a coordinate speed alone cannot answer\textemdash{}whether the metric is Lorentzian, whether the carrier is locally timelike, and whether distant clocks register an arrival advantage. Its conclusions are conditional on a prescribed geometry, rather than on an independently established matter source.

Use signature $(-,+,+,+)$ and coordinates $(t,x,y,z)$, with $t$ a physical time coordinate, the spatial coordinates measured in length units, and $c$ the local photon speed. Let the center position $x_s(t)$ be smooth and nondecreasing, with $x_s(t)=0$ for $t\leq 0$ and $x_s(t)=L$ for $t\geq T$, where $L>cT$. Define $v(t)=dx_s/dt$ and $r^2=[x-x_s(t)]^2+y^2+z^2$, and choose a smooth dimensionless radial profile satisfying $0\leq f(r)\leq 1$, equal to unity near the center and vanishing for $r\geq R$. The target metric is

$ds^2=-c^2dt^2+[dx-v(t)f(r)dt]^2+dy^2+dz^2.$

The shifted spatial differential changes the relation between coordinate displacement and local motion without changing the metric signature. In particular, the coframe $(c\,dt,\,dx-vf\,dt,\,dy,\,dz)$ is invertible and puts the line element into explicitly Lorentzian form; the coordinate determinant is $\det(g_{\mu\nu})=-c^2$. Thus even an interval with $v(t)>c$ does not make the metric degenerate or change its signature. Because the center is stationary before and after the prescribed motion, and because the profile has compact support, the metric is exactly Minkowski for $t\leq 0$, for $t\geq T$, and outside the moving region $r<R$.

The central carrier follows $x=x_s(t)$ and $y=z=0$, where the profile is constant at $f=1$. Along this trajectory, $dx=v(t)dt$, so the shifted spatial differential vanishes and

$ds^2=-c^2dt^2,\qquad d\tau=dt.$

The carrier is therefore timelike throughout the transit, and its elapsed proper time between departure and arrival is $T$. A large displacement relative to the exterior coordinates does not, by itself, represent locally superluminal motion.

The same distinction follows directly from the local causal cone. Relative to observers normal to the constant-time slices, define the infinitesimal spatial components $d\ell_x=dx-vf\,dt$, $d\ell_y=dy$, and $d\ell_z=dz$; causal displacements satisfy

$d\ell_x^2+d\ell_y^2+d\ell_z^2\leq c^2dt^2.$

The central carrier has vanishing spatial components in this frame, even when its exterior coordinate velocity exceeds $c$. The shortcut consequently does not grant a massive carrier permission to cross its local photon cone; it changes the spacetime geometry used to relate local motion to distant endpoints.

To define those endpoints operationally, place stationary station A at $x=-D$ and station B at $x=L+D$, both on the axis, with $D>R$. Monotonicity of the center trajectory keeps both stations outside the modified region for the entire pulse, so their proper-time increments equal $dt$. Synchronize their clocks in the common Minkowski past. At $t=-D/c$, station A emits a light signal that reaches the central carrier at $(t,x)=(0,0)$, entirely through the flat pre-pulse region.

The carrier retains the signal information until it reaches $(t,x)=(T,L)$ and then emits a light signal toward station B. This final light segment lies entirely in the flat post-pulse region and reaches B at $t=T+D/c$. The elapsed time obtained by subtracting the synchronized exterior-clock readings is therefore

$\Delta t_{\mathrm{shortcut}}=T+\frac{2D}{c}<\frac{L+2D}{c}=\Delta t_{\mathrm{Minkowski}},$

where the strict inequality is precisely the imposed condition $L>cT$. The reference on the right is the photon transit between the same stationary endpoints in unmodified Minkowski spacetime, not the fastest null transit in the modified geometry.

This comparison uses two geometries with matched exterior clocks, rather than a timelike carrier racing a null signal through the same metric. The stations, carrier, and relays are idealized test systems whose backreaction is neglected, and the geometric schedule is prearranged independently of the newly transmitted information. The protocol establishes an arrival advantage within that assigned geometry; it does not establish that station A can create the geometry on demand or that an admissible dynamical source generates it.

The assigned spacetime also admits a direct causal-loop test that is independent of the transit comparison. Its inverse metric obeys

$g^{tt}=-\frac{1}{c^2},$

so the gradient of $t$ is timelike everywhere on the stated coordinate domain. Choosing the future orientation in which $dt$ is positive on future-directed causal vectors makes $t$ strictly increase along every nontrivial future-directed causal curve, excluding closed causal curves and supplying the usual stable-causality criterion. This result applies to this spacetime, not automatically to assemblies of relatively moving shortcuts or additional transport channels, whose combined causal structure would require a separate analysis.

The finite-duration construction therefore defeats the purely kinematic inference that an effective shortcut must entail either locally superluminal carrier motion or a causal loop. What remains at issue is whether a physical source and its evolution can realize this operational geometry, not whether its exterior coordinate speed alone contradicts its local light cones.

\section{The Einstein source: negative energy and null stress}
The Einstein source calculation turns the prescribed shortcut geometry into a precise material requirement: negative energy measured by the slice-normal observers and, independently, strict negative null stress in the wall. Consider the constant-speed geometry $ds^2=-(dx^0)^2+[dx-u\,dx^0]^2+dy^2+dz^2$, with $x^0=ct$, $u=\beta f(r)$, constant $\beta=v/c$, and $r^2=(x-\beta x^0)^2+y^2+z^2$. The dimensionless profile is smooth, equals one for $r\leq R$, and vanishes for $r\geq R+\Delta$, where $R,\Delta>0$ are lengths. Throughout this calculation, assume signature $(-,+,+,+)$, zero cosmological constant, and Einstein\textquoteright{}s equation $G_{ab}=8\pi G T_{ab}/c^4$.

The metric has unit lapse, Euclidean spatial metric $h_{ij}=\delta_{ij}$, and shift $B^i=(-u,0,0)$. An orthonormal frame consists of the future-directed unit normal $n=\partial_0+u\partial_x$ and spatial vectors $e_i=\partial_i$. Define the measured projections by $\varepsilon=T(n,n)$, $j_i=-T(n,e_i)$, and $S_{ij}=T(e_i,e_j)$; consequently, the covariant frame components satisfy $T_{0i}=-j_i$. These projections all have units $\mathrm{J\,m^{-3}}$, and the momentum-flux sign is essential when forming null contractions.

Adopt the extrinsic-curvature convention $K_{ij}=-(\mathcal L_n h)_{ij}/2$. Since the spatial metric is time independent, $K_{ij}=(\partial_i B_j+\partial_j B_i)/2$, so its nonzero independent components are $K_{xx}=-\beta f_x$, $K_{xy}=-\beta f_y/2$, and $K_{xz}=-\beta f_z/2$, where subscripts on the profile denote spatial partial derivatives. Thus $K^2=\beta^2 f_x^2$, whereas $K_{ij}K^{ij}=\beta^2 f_x^2+\beta^2(f_y^2+f_z^2)/2$. The longitudinal derivative contributes equally to these two quadratic expressions; only the transverse derivatives survive their difference.

The Hamiltonian constraint is ${}^{(3)}R+K^2-K_{ij}K^{ij}=16\pi G\varepsilon/c^4$. Because the slices are intrinsically flat, it gives $\varepsilon=-c^4\beta^2(f_y^2+f_z^2)/(32\pi G)$. For nonzero $\beta$, the energy density is therefore strictly negative wherever the transverse wall gradient is nonzero, violating the weak energy condition and hence the dominant energy condition. For a spherical profile, $f_y^2+f_z^2=[f'(r)]^2\sin^2\theta$, with $\theta$ measured from the translation axis, so every nontrivial smooth wall contains an open negative-density region.

The angular weighting also fixes the coefficient in the integrated requirement. The proper spatial volume is $d^3x$, and $\int\sin^2\theta\,d\Omega=8\pi/3$, giving the normal-frame slice integral $E_n\equiv\int\varepsilon\,d^3x=-[c^4\beta^2/(12G)]\int_R^{R+\Delta}r^2[f'(r)]^2\,dr$. The radial weight matters: a bound based only on the coordinate thickness would omit the spherical wall\textquoteright{}s area dependence.

This negative volume integral is not the asymptotic gravitational energy or a demonstrated propulsion expenditure. The standard Arnowitt\textendash{}Deser\textendash{}Misner energy surface integral vanishes because the spatial metric is exactly Euclidean, and the compact support of the extrinsic curvature likewise gives zero asymptotic momentum. Thus negative $E_n$ must not be identified with negative total mass, and its magnitude is not an established fuel requirement. There is no contradiction with a positive-energy theorem whose usual dominant-energy premise is violated by the required source.

A profile-independent lower bound follows from the endpoint change, without assuming monotonicity. Since $\int_R^{R+\Delta}f'(r)\,dr=-1$, weighted Cauchy\textendash{}Schwarz gives $1\leq[\int_R^{R+\Delta}r^2[f'(r)]^2\,dr][\int_R^{R+\Delta}r^{-2}\,dr]$. The second factor is $\Delta/[R(R+\Delta)]$, and therefore $|E_n|\geq[c^4\beta^2/(12G)]R(R+\Delta)/\Delta$. Equality would require the radial derivative to be proportional to $r^{-2}$ throughout the wall, which cannot smoothly join both constant-profile regions.

For a thin wall, $\Delta\ll R$, the leading lower-bound scale is $c^4\beta^2R^2/(12G\Delta)$: reducing the transition thickness increases the required magnitude rather than removing the negative energy. Since $c^4/G$ has units $\mathrm{J\,m^{-1}}$, the bound has energy units. Its scope is the stated spherical, unit-lapse, flat-slice ansatz with $R>0$, not arbitrary shortcut geometries or a universal cost of effective superluminal travel. The normal density and its slice integral vanish when $\beta=0$ or the profile is constant.

Negative normal energy alone does not establish violation of the null energy condition; the spatial stresses and momentum flux must also be included. To obtain the relevant components, write $C=c^4/(8\pi G)$, $Q=u_y^2+u_z^2$, and $L=u_{yy}+u_{zz}$. With the stated signs, the momentum constraint is $j_i/C=\partial_jK^j{}_i-\partial_iK$, yielding $j_x=-CL/2$, while the Hamiltonian result becomes $\varepsilon=-CQ/4$. Both $Q$ and $L$ have units $\mathrm{m^{-2}}$.

The longitudinal stress follows from the spatial Einstein evolution equation rather than an assumed equation of state. With $S=\delta^{ij}S_{ij}$, this equation reads $(\partial_0-\mathcal L_B)K_{ij}=KK_{ij}-2K_{ik}K^k{}_j-[S_{ij}-\delta_{ij}(S-\varepsilon)/2]/C$. Define $a=u_x$ and let $D=n^a\partial_a$ act on scalars; for the translating profile, $D=(u-\beta)\partial_x$. If $P_{ij}$ denotes the geometric combination $KK_{ij}-2K_{ik}K^k{}_j-(\partial_0-\mathcal L_B)K_{ij}$, its longitudinal component and trace are $P_{xx}=Da+a^2-Q/2$ and $P=Da+a^2$.

Tracing and then inverting the evolution equation gives $S/C=3\varepsilon/C-2P$ and $S_{ij}/C=P_{ij}+\delta_{ij}(\varepsilon/C-P)$. Consequently, $S_{xx}=C[(Da+a^2-Q/2)-Q/4-(Da+a^2)]=-3CQ/4$. The longitudinal derivative and evolution terms cancel in this stress projection, leaving a negative contribution fixed by the same transverse-gradient quantity that determines the normal density.

For the future-directed null vectors $k_\pm=n\pm e_x$, the measured source components now give $T(k_\pm,k_\pm)=\varepsilon+S_{xx}\mp2j_x=C(-Q\pm L)$. Equivalently, $T(k_\pm,k_\pm)=[c^4/(8\pi G)][-\beta^2(f_y^2+f_z^2)\pm\beta(f_{yy}+f_{zz})]$. The second-derivative flux term changes sign between the two directions, but the transverse-gradient term is negative in both.

The same cancellation is visible directly in curvature. Using $R(X,Y)Z=\nabla_X\nabla_YZ-\nabla_Y\nabla_XZ-\nabla_{[X,Y]}Z$, the orthonormal projections are $R_{nn}=-Da-a^2-Q/2$, $R_{xx}=Da+a^2-Q/2$, and $R_{nx}=L/2$. Since contraction with a null vector removes the scalar-curvature term in Einstein\textquoteright{}s equation, these expressions reproduce $T(k_\pm,k_\pm)=C(-Q\pm L)$. This supplies a second analytic route to the null result, distinct from inferring it from the negative density.

The two contractions sum to $-2CQ$, so at least one is strictly negative wherever $Q>0$. Their minimum is $C(-Q-|L|)$; if $Q=0$ but $L\neq0$, one direction still violates the null energy condition, whereas $Q=L=0$ leaves other null directions undecided. Every nontrivial smooth spherical wall with $\beta\neq0$ has an open region with $Q>0$, independently of monotonicity and without requiring $|\beta|>1$. The source obstruction is therefore a property of this metric family even below the nominal superluminal threshold, not an effect that first appears at that threshold.

These are pointwise statements, not an averaged null-energy theorem. An averaged condition requires a specified complete, affinely parametrized null geodesic, its tangent normalization, and convergence of the stress integral; the locally defined pair $k_\pm$ cannot simply be integrated as if it supplied those data. Likewise, defining $T_{ab}=c^4G_{ab}/(8\pi G)$ ensures $\nabla^aT_{ab}=0$ by the contracted Bianchi identity, but supplies neither a material constitutive law nor preparable data with stable coupled evolution.

The calculation therefore excludes everywhere-null-energy-respecting matter as the source of this exact target and makes reproduction of its negative null stress\textemdash{}not merely its geometry\textemdash{}the necessary next condition for any proposed transport mechanism.

\section{Formation from complete data and the canonical-scalar counterexample}
Formation is a constraint on complete dynamical data, not merely on the geometry prescribed for a later transit. Under autonomous classical equations with a unique Cauchy development, complete vacuum data cannot evolve into a nontrivial curvature pulse. A canonical scalar with a sign-changing potential provides a substantive test of whether initially hidden, stress-free matter can evade this obstruction: negative energy is algebraically available, but the admissible stress-free data and their subsequent evolution still fail to produce the required source.

The vacuum obstruction requires three explicit hypotheses. The autonomous Einstein\textendash{}matter theory must admit Minkowski spacetime with matter vacuum $\Psi_0$; the initial surface must be a complete Cauchy surface $\Sigma\cong\mathbb{R}^3$ carrying $h_{ij}=\delta_{ij}$, $K_{ij}=0$, $\Psi=\Psi_0$, and every independent matter momentum at its vacuum value; and the smooth maximal globally hyperbolic development must be unique up to gauge and diffeomorphism. Minkowski vacuum is already a development of these data, so uniqueness excludes an inequivalent later pulse within that development.

Incoming preparation must therefore be represented in the complete initial data, or enter through a boundary if the chosen surface was not actually a Cauchy surface.

Flat gravitational data alone do not establish the matter-vacuum premise. The gravitational constraints fix particular projections of the stress tensor, whereas an exactly Minkowski spacetime region fixes the entire tensor through Einstein's equation. An external actuator likewise cannot be omitted from the source description: its fields, stress, and preparation data belong to the enlarged system. These distinctions separate a genuine vacuum-formation obstruction from the stronger, generally invalid inference that initially flat geometry contains no dynamical matter.

For an explicit classical test, adopt signature $(-,+,+,+)$, natural units $c=\hbar=1$, ordinary Einstein gravity with zero cosmological constant, and $\kappa=8\pi G$. The minimally coupled real scalar is specified by

$S=\int d^4x\,\sqrt{-g}\left[\frac{R}{2\kappa}-\frac12 g^{\mu\nu}\partial_\mu\phi\,\partial_\nu\phi-V(\phi)\right],\qquad V(\phi)=\frac{\lambda}{a^2}\phi^2(\phi^2-a^2)(\phi^2-4a^2),\quad \lambda>0,\quad a>0.$

Here $\phi$ and $a$ have energy units, $\lambda$ is dimensionless, $V$ has energy-density units, and $\kappa$ has units of inverse energy squared. This polynomial defines a classical source model, not a demonstrated quantum-complete theory.

The potential is bounded below but negative for $a<|\phi|<2a$, so rejecting the source cannot rest on an assumed nonnegative potential. Its vacuum at $\phi=0$ satisfies $V(0)=V'(0)=0$ and $V''(0)=8\lambda a^2>0$. Variation gives

$G_{\mu\nu}=\kappa T_{\mu\nu},\qquad \Box\phi-V'(\phi)=0,\qquad T_{\mu\nu}=\partial_\mu\phi\,\partial_\nu\phi-g_{\mu\nu}\left[\frac12(\partial\phi)^2+V(\phi)\right],$

where $(\partial\phi)^2=g^{\alpha\beta}\partial_\alpha\phi\,\partial_\beta\phi$. The analysis concerns smooth solutions in their unique regular globally hyperbolic development; it assumes no extension through a singularity.

Let $n$ denote the unit future normal to an initial slice, $h_{ij}$ its induced metric, $\pi=n^\mu\partial_\mu\phi$ the scalar normal derivative, and $q_i=D_i\phi$ its spatial gradient, with $q^2=h^{ij}q_iq_j$. The normal energy density, momentum density, and spatial stress are

$\rho=\frac12(\pi^2+q^2)+V,\qquad j_i=-\pi q_i,\qquad S_{ij}=q_iq_j+h_{ij}\left[\frac12(\pi^2-q^2)-V\right].$

Complete Einstein\textendash{}scalar data comprise $(h,K,\phi,\pi)$ subject to the gravitational constraints. In particular, $h_{ij}=\delta_{ij}$ and $K_{ij}=0$ require $\rho=0$ and $j_i=0$, but do not by themselves require $S_{ij}=0$.

Requiring the full tensor to vanish gives a much sharper classification than requiring only its constraint projections to vanish. For every unit spatial vector $e$, the scalar tensor obeys

$T(n,n)+T(e,e)=\pi^2+(e\cdot q)^2.$

If $T_{\mu\nu}=0$, this nonnegative expression vanishes for every $e$, forcing $\pi=0$ and $q_i=0$; the remaining condition $\rho=0$ then gives $V=0$. Conversely, these three conditions make every component of the stress tensor vanish.

On a connected slice, all smooth stress-free scalar data therefore have spatially constant field value $\phi=C$, with $C\in\{0,\pm a,\pm2a\}$. This result uses no assumption that $V\geq0$ and leaves no smooth localized stress-free profile. The asymptotic matter condition $\phi\to0$ retains only $C=0$, but relaxing that condition admits four nonvacuum, stress-free constants. They are genuine alternatives to complete matter-vacuum data, even though their initial stress tensors vanish.

The nonzero roots are not stationary solutions, because $V'(\pm a)=\mp6\lambda a^3$ and $V'(\pm2a)=\pm48\lambda a^3$. Thus the complete data $(\delta,0,C,0)$ satisfy both gravitational constraints and have $T_{\mu\nu}=0$ initially, yet the scalar equation supplies a nonzero acceleration whenever $C\ne0$. These data are a concrete counterexample to identifying an initially vanishing Einstein tensor with complete vacuum data. They do not contradict vacuum-data uniqueness, because they omit its matter-vacuum premise.

Their evolution can be determined without assuming that the initially hidden field becomes a localized wall. Translations and rotations preserve $(\delta,0,C,0)$, and uniqueness preserves those symmetries in the development. Locally, the solution consequently has spatially flat Friedmann\textendash{}Lemaître\textendash{}Robertson\textendash{}Walker form

$ds^2=-d\tau^2+b(\tau)^2d\mathbf{x}^2,\qquad \phi=\phi(\tau),$

where $\tau$ is proper time, $b$ is the dimensionless scale factor, and $H=\dot b/b$, with a dot denoting differentiation with respect to $\tau$. The initial conditions are $b(0)=1$, $H(0)=0$, $\phi(0)=C$, and $\dot\phi(0)=0$.

The homogeneous coupled equations are

$3H^2=\kappa\rho,\qquad \dot H=-\frac{\kappa}{2}\dot\phi^2,\qquad \ddot\phi+3H\dot\phi+V'=0,\qquad \rho=\frac12\dot\phi^2+V.$

The scalar equation gives the conservation law $\dot\rho=-3H\dot\phi^2$, and hence $d(3H^2-\kappa\rho)/d\tau=0$. This verifies propagation of the Hamiltonian constraint without dividing by the initially vanishing $H$. On every regular portion of the development, it follows that $\rho=3H^2/\kappa\geq0$.

This nonnegative energy is not an artifact of choosing the homogeneous normal. With homogeneous pressure $p=\dot\phi^2/2-V$, one has $\rho+p=\dot\phi^2$, and any unit timelike observer $w$ with Lorentz factor $\gamma$ relative to that normal measures

$T(w,w)=(\rho+p)\gamma^2-p=\rho+(\gamma^2-1)\dot\phi^2\geq0.$

Thus every stress-free flat-slice development classified above satisfies the weak energy condition throughout its regular evolution, even when the field enters the negative-potential region. Spatial homogeneity also prevents these data from generating a compact moving wall.

A local expansion shows explicitly why initially vanishing stress does not imply a flat time interval. Choose $C=a$ and define $A=-V'(a)=6\lambda a^3$; for $|V''(a)|\tau^2\ll1$ and $\kappa A^2\tau^4\ll1$, the smooth equations give

$\phi=a+\frac{A}{2}\tau^2+O(\tau^4),\qquad H=-\frac{\kappa A^2}{6}\tau^3+O(\tau^5),\qquad \rho=\frac{\kappa A^4}{12}\tau^6+O(\tau^8),\qquad p=A^2\tau^2+O(\tau^4).$

Kinetic and potential energy cancel at quadratic order in proper time, but pressure does not, and the Ricci scalar becomes $R=\kappa(\rho-3p)=-3\kappa A^2\tau^2+O(\tau^4)$. This is a short-time analytical expansion, not a stability result; it demonstrates curvature on both sides of the initially stress-free slice rather than an exactly flat prehistory followed by a pulse.

An exactly Minkowski open prehistory eliminates these nonstationary alternatives more directly. Einstein's equation requires $T_{\mu\nu}=0$ throughout that region, so the stress-free classification forces the scalar to be spacetime-constant there; its field equation additionally requires $V'=0$. Among the zero-potential roots, only $C=0$ satisfies both requirements. A complete slice in that prehistory therefore carries precisely the vacuum data $(\delta,0,0,0)$, and uniqueness prohibits subsequent pulse formation, regardless of a prescribed smooth switching schedule.

There is also a source exclusion stronger than the prehistory argument. For every metric-null vector $k$, the canonical scalar obeys $T(k,k)=(k^\mu\partial_\mu\phi)^2\geq0$, independently of the sign of its potential, whereas the exact shortcut target requires a strictly negative null contraction somewhere in its moving wall. Consequently, even inhomogeneous preparation in the negative-potential region cannot make this unchanged scalar source the exact target. The formation analysis therefore changes the task from finding a better preparation of this scalar to changing the source theory or the target geometry; it excludes this source\textendash{}target combination, not effective superluminal travel in every framework.

\section{A local stability gate for nonlinear scalar sources}
Nonlinear first-derivative scalar dynamics do not remove the target\textquoteright{}s source obstruction if the scalar must retain positive high-frequency energy and a regular wave equation on a common matter\textendash{}gravity Cauchy slicing. The relevant exclusion concerns a single, smooth, minimally coupled real scalar with Lagrangian $P(X,\phi)$ in ordinary positive-sign Einstein gravity; any additional sources are assumed to satisfy the null energy condition. The decisive conflict is local: the sign required for strictly negative null stress also fixes a negative transverse contribution to the scalar perturbation energy.

For this analysis only, adopt natural units $c=\hbar=1$, signature $(-,+,+,+)$, and $\kappa=8\pi G$. The action and kinetic invariant are $S=\int d^4x\,\sqrt{-g}\,[R/(2\kappa)+P(X,\phi)]$ and $X=-g^{\mu\nu}\partial_\mu\phi\partial_\nu\phi/2$, with $P$ smooth and its derivatives finite on the background. Taking $\phi$ to have energy units makes $X$ and $P$ have units of energy to the fourth power, $P_X$ dimensionless, and $P_{XX}$ have units of inverse energy to the fourth power.

The necessary stability gate requires a physical scalar mode with a nondegenerate Lorentzian principal symbol, a common spacelike initial surface with gravity, and positive high-frequency kinetic and spatial-gradient energies relative to the gravitational sector; it is not a sufficient test of nonlinear stability or quantum consistency.

The source requirement entering this test is strict negative null stress, not merely negative normal energy. For the target metric $ds^2=-dt^2+(dx-b\,dt)^2+dy^2+dz^2$, define $n=\partial_t+b\partial_x$, $e_x=\partial_x$, $k_\pm=n\pm e_x$, and $Q=(\partial_y b)^2+(\partial_z b)^2$. The curvature contraction gives $T(k_+,k_+)+T(k_-,k_-)=-2Q/\kappa$, including time-dependent profiles, so at least one null contraction is strictly negative wherever $Q>0$. This is the geometric requirement that the scalar must meet, independently of how its initial data are prepared.

Write $u_\mu=\partial_\mu\phi$, $A=P_X$, and $B=P_{XX}$. Variation of the action gives $T_{\mu\nu}=A u_\mu u_\nu+P g_{\mu\nu}$ and the scalar equation $\mathcal E\equiv\nabla_\mu(A\nabla^\mu\phi)+P_\phi=0$. For every metric-null vector $k^\mu$, the contraction is therefore $T_{\mu\nu}k^\mu k^\nu=A(k^\mu u_\mu)^2$: strict negativity requires $A<0$ and a nonzero contraction, while the term proportional to the metric cannot change its sign. The identity $\nabla_\mu T^{\mu\nu}=\mathcal E\nabla^\nu\phi$ also ensures covariant stress conservation on scalar solutions.

The same derivatives of $P$ determine the scalar\textquoteright{}s propagation and perturbation energy. Expanding about a smooth background solution as $\phi=\bar\phi+\pi$, and integrating the mixed field\textendash{}derivative term by parts, gives the fixed-metric quadratic Lagrangian $\mathcal L_2=-Z^{\mu\nu}\nabla_\mu\pi\nabla_\nu\pi/2+U\pi^2/2$, where $Z^{\mu\nu}=A g^{\mu\nu}-B u^\mu u^\nu$ and $U=P_{\phi\phi}+\nabla_\mu(P_{X\phi}u^\mu)$, evaluated on the background. Here $Z^{\mu\nu}$ is also the principal coefficient of the scalar equation, and $U$ has units of energy squared.

Freezing the coefficients supplies a necessary short-wavelength test: lower-derivative terms and background variation cannot reverse its leading derivative signs in a controlled high-frequency regime.

To expose those signs without assuming a particular background gradient, choose a local metric-orthonormal frame adapted to an arbitrary metric-timelike normal. Let $q=\partial_t\bar\phi$, $s_i=\partial_i\bar\phi$, and $K=A+Bq^2$, where the dot below denotes differentiation along the chosen local time direction. For $K\ne0$, the conjugate momentum is $p=K\dot\pi-Bq(\mathbf s\cdot\boldsymbol\nabla\pi)$, and the quadratic Hamiltonian density is $\mathcal H_2=[p+Bq(\mathbf s\cdot\boldsymbol\nabla\pi)]^2/(2K)+[A|\boldsymbol\nabla\pi|^2-B(\mathbf s\cdot\boldsymbol\nabla\pi)^2]/2-U\pi^2/2$. Thus changing the slicing changes the kinetic coefficient and the mixed term, but it does not eliminate the transverse coefficient $A$.

Choose a perturbation with spatial gradient $\mathbf d=\boldsymbol\nabla\pi$ perpendicular to $\mathbf s$. Its derivative Hamiltonian reduces to $\mathcal H_{2,\mathrm{der}}=p^2/(2K)+A|\mathbf d|^2/2$, so positive momentum and transverse-gradient energies require both $K>0$ and $A>0$. If the required $A<0$ holds, then $K<0$ gives negative momentum energy, whereas $K>0$ leaves negative transverse-gradient energy; $K=0$ makes the chosen surface characteristic for the scalar. On the formal transverse Fourier modes with wavevector $\mathbf k\perp\mathbf s$, the frozen equation gives $\omega^2=(A|\mathbf k|^2-U)/K$, which is negative at sufficiently large $|\mathbf k|$ when $K>0$, with growth rate approaching $|\mathbf k|\sqrt{-A/K}$.

The interpretation of this exponential behavior depends on whether the proposed initial surface is actually a scalar Cauchy surface. For $A<0$, the two spatial directions transverse to $\mathbf s$ already give two negative directions of $Z^{\mu\nu}$; if the symbol is Lorentzian, it must consequently have three negative eigenvalues and one positive eigenvalue. A scalar Cauchy normal then obeys $Z^{\mu\nu}n_\mu n_\nu>0$, which here means $K<0$, exposing a ghost\textemdash{}a physical mode with negative kinetic energy relative to the ordinary gravitational sector.

Formal exponential solutions found with $K>0$ can instead diagnose an inadmissible time slicing, rather than an instability evolved from valid scalar Cauchy data; neither interpretation restores the required positive-energy common evolution.

The gradient classifications make clear why hyperbolicity alone is insufficient. For a timelike scalar gradient, $X>0$, use its rest frame and define $F=A+2XB$; then $Z^{\mu\nu}=\operatorname{diag}(-F,A,A,A)$ and $\mathcal H_2=p^2/(2F)+A|\boldsymbol\nabla\pi|^2/2-U\pi^2/2$. Positive high-frequency energy requires $F>0$ and $A>0$, with sound speed squared $c_s^2=A/F$. For the required $A<0$, choosing $F>0$ produces a definite rather than Lorentzian symbol and negative squared transverse frequencies, while $F<0$ restores a Lorentzian symbol only with negative kinetic energy. Thus the hyperbolicity condition $AF>0$ permits an overall wrong-sign branch and cannot substitute for the energy test.

For a spacelike scalar gradient, $X<0$, align the local spatial axis with $u_\mu=(0,\sqrt{-2X},0,0)$. The principal symbol becomes $Z^{\mu\nu}=\operatorname{diag}(-A,F,A,A)$, and the derivative Hamiltonian is $p^2/(2A)+F(\partial_x\pi)^2/2+A[(\partial_y\pi)^2+(\partial_z\pi)^2]/2$. Positive energy again requires $A>0$ and $F>0$; altering $B$ changes the longitudinal coefficient but cannot repair the negative kinetic coefficient when $A<0$. With $AF<0$ the symbol has two eigenvalues of each sign rather than Lorentzian signature, while $A<0$ and $F<0$ give a Lorentzian but wrong-sign wave operator.

A nonzero null gradient does not provide an exception. In a local frame with $u_\mu=(q,q,0,0)$ and $X=0$, the determinant is $\det Z=-A^4$, so finite $B$ and $A\ne0$ retain a Lorentzian symbol, but the general transverse Hamiltonian still fails the positive-energy test when $A<0$. For $A>0$, a local boost can make $|q|$ sufficiently small to obtain positive kinetic and spatial-gradient coefficients, but the null stress is then nonnegative. At zero background gradient, $u_\mu=0$, all null stresses vanish and $Z^{\mu\nu}=A g^{\mu\nu}$; negative $A$ is simply the wrong-sign scalar, while $A=0$ removes the quadratic derivative term.

Degenerate boundaries must be separated from regular propagating scalars rather than declared universally inconsistent. For non-null gradients, $\det Z=-A^3F$: at $A=0$ the null stress vanishes and the principal operator loses rank, while at $F=0$ it also becomes degenerate. In the timelike-gradient rest frame, $A=0$ leaves only a quadratic time derivative and no spatial restoring term, whereas $F=0$ leaves a spatial constraint rather than a scalar wave equation; for spacelike gradients, the corresponding limits remove either the time derivative or the longitudinal principal derivative. For a nonzero null gradient with $A=0$, $Z^{\mu\nu}=-B u^\mu u^\nu$ has rank one when $B\ne0$ and rank zero otherwise.

None of these cases supplies the stipulated nondegenerate wave cone.

A concrete constrained boundary is $P=-\mu^2\sqrt{2X}-V(\phi)$ on $X>0$, for which $A=-\mu^2/\sqrt{2X}<0$ but $F=0$. It can violate the null energy condition, yet its scalar principal part is a spatial constraint in the gradient rest frame, so it lies outside the regular propagating-mode criterion. Its predictivity and stability require a separate coupled constraint and elliptic\textendash{}hyperbolic analysis, including boundary conditions; failure of the present gate is not a proof that every such formulation is incoherent. Higher-derivative terms likewise change the specified action and require a new test, while smooth passage through $A=0$ cannot remove the sign obstruction within any open region of strict negative null stress.

For regular minimal coupling, the scalar principal-part test remains a necessary local gate for the coupled Einstein\textendash{}scalar system in the controlled high-frequency regime. Applied as an effective field theory, it requires a nonempty wavelength window short compared with background variation but still below the theory\textquoteright{}s cutoff; without that window, the conclusion is restricted to the exact classical model rather than unknown ultraviolet dynamics. The result therefore changes the source search, not the kinematic definition of a shortcut: changing preparation cannot make a regular positive-energy single $P(X,\phi)$ scalar support this exact target.

Further progress requires a changed action, a changed target, or an explicitly analyzed constrained formulation, rather than treating nonlinear dependence on the first-derivative invariant as a stable source mechanism.

\section{Axial completeness, tidal curvature, and averaged energy}
Affine propagation imposes a distinct test on a prescribed shortcut geometry: local negative null stress, a negative integral along a complete ray, and a finite-affine tidal endpoint are different findings. Consider the eternal constant-speed metric $ds^2=-(dx^0)^2+[dx-\beta f(r)dx^0]^2+dy^2+dz^2$ on $M=\mathbb{R}^4$, where $x^0=ct_{\mathrm{phys}}$, $\xi=x-\beta x^0$, and $r^2=\xi^2+y^2+z^2$. The signature is $(-,+,+,+)$, and the smooth spherical profile equals one for $r\leq R$ and zero for $r\geq R+\Delta$, with $R,\Delta>0$. The analysis is classical, with zero cosmological constant and required stress tensor $T_{ab}=C G_{ab}$, where $C=c^4/(8\pi G)$; quantum and test-ray backreaction are not included.

The axial equations make completeness accessible before any averaged-energy interpretation is assigned. Write $a(\xi)=1-\beta+\beta f(|\xi|)$ and $w=\beta(1-f)=1-a$, so the comoving metric is $ds^2=-(dx^0)^2+(d\xi+w\,dx^0)^2+dy^2+dz^2$. Reflection symmetry preserves axial motion, and its future-directed forward null tangent is $\ell=\omega(\partial_0+a\partial_\xi)$, corresponding to $\omega(n+e_x)$ in the unit-normal frame, with $\omega=dx^0/d\lambda>0$. Here $\lambda$ has units of length, $\omega$ is dimensionless, and $\partial_0$ denotes differentiation at fixed comoving coordinates.

The stationary Killing field $K=\partial_0$ gives $E=-g(K,\ell)=\omega a$, while the affine equations give $d\omega/d\lambda=-\omega^2\,da/d\xi$ and $d\xi/d\lambda=E$.

For the restricted subluminal case $0<\beta<1$, assume additionally that the nontrivial profile is monotonically decreasing through its wall and satisfies $0\leq f\leq1$. Then $1-\beta\leq a\leq1$, so $a$ never vanishes. Exterior normalization $\omega=1$ fixes $E=1-\beta$, and consequently $1-\beta\leq\omega\leq1$. Since $\xi$ evolves affinely with a nonzero constant derivative and the frequency remains bounded, the forward axial rays extend over a two-sided infinite affine domain; their stress contraction is supported only in the wall crossings.

Completeness permits the averaged null-energy condition to be tested using $I_\gamma=\int_{-\infty}^{+\infty}T_{ab}\ell^a\ell^b\,d\lambda$. For this subluminal family, the axial contraction and the two symmetric wall crossings give $I_+=4C\beta(1-\beta)\int_R^{R+\Delta}f'(r)\big/\{r[1-\beta+\beta f(r)]^2\}\,dr<0$, where $f'$ denotes the radial derivative. The integral is finite because its denominator is bounded away from zero, $R>0$, and the profile derivative has compact support; its strict negativity follows from the nontrivial monotone decrease.

This is a complete-geodesic averaged-null-energy violation under the stated normalization and profile assumptions, not a result that can be transferred to the superluminal regime by substituting a larger value of $\beta$.

For the superluminal analysis, fix $\beta>1$ and an explicit smooth profile rather than assuming that every compact wall has the same endpoint structure. Define $\eta(s)=\exp(-1/s)$ for $s>0$ and $\eta(s)=0$ for $s\leq0$, and set $s=(r-R)/\Delta$ and $f(r)=\eta(1-s)/[\eta(s)+\eta(1-s)]$. This profile is $C^\infty$, is constant in the core and exterior, and decreases strictly within the wall. The constant core also ensures smoothness of the radial construction at the origin.

There is exactly one $h\in(R,R+\Delta)$ satisfying $f(h)=1-1/\beta$, so the forward axial coefficient vanishes precisely at $\xi=-h$ and $\xi=+h$. In the wall, $f'(r)=-f(1-f)[s^{-2}+(1-s)^{-2}]/\Delta<0$; therefore, with $\kappa=-\beta f'(h)>0$, the zero slopes are $a'(-h)=\kappa$ and $a'(+h)=-\kappa$. Thus $a$ is positive between the zeros and negative outside them, and both zeros are simple. These statements identify axial characteristic separatrices without requiring the full spherical surface $r=h$ to be a global event horizon.

For $E\neq0$, the first integral $d\xi/d\lambda=E$ divides the forward rays into three sectors, with the sign of $E$ fixed by $\omega>0$. In $\xi<-h$, one has $E<0$: the future runs to $\xi=-\infty$ at infinite affine parameter, whereas the past approaches $\xi=-h$ at finite affine parameter. In $-h<\xi<h$, one has $E>0$, and the ray travels from $-h$ to $+h$ over the finite affine interval $2h/E$. In $\xi>h$, one again has $E<0$: the past extends to $\xi=+\infty$ at infinite affine parameter, while the future approaches $\xi=+h$ at finite affine parameter. The exterior ends are genuinely affine-complete because $a=1-\beta$ is constant there.

The finite-affine boundaries are not finite-coordinate-time points at which these rays can simply cross into another sector. Near a zero $\xi_h$, let $q=a'(\xi_h)=\pm\kappa$; integrating $dx^0/d\xi=1/a$ gives $x^0=q^{-1}\ln(|\xi-\xi_h|/\Delta)+O(1)$. Every approach to $-h$ therefore occurs as $x^0\to-\infty$, and every approach to $+h$ occurs as $x^0\to+\infty$. The resulting rays are maximal in the specified manifold because their finite-affine limits have no finite-coordinate-time endpoint in $M$.

The zero-energy generators must be treated separately because division by $a$ would omit them. A nonzero future null tangent with $E=0$ lies at one of the zeros, with $\xi=\xi_h$ and $y=z=0$ for all real coordinate times. Its frequency satisfies $\omega=1/[q(\lambda-\lambda_c)]$, and integration gives $\lambda-\lambda_c=\exp[q(x^0-x^0_0)]/(q\omega_0)$ for $\omega_0>0$. At $\xi=-h$, the affine domain is $(\lambda_c,\infty)$, so the generator is past incomplete and future complete; at $\xi=+h$, the domain is $(-\infty,\lambda_c)$, giving the opposite classification.

Together with the three nonzero-energy sectors, these generators establish that every maximal forward axial ray is incomplete in at least one affine direction, even when its orbit covers every finite coordinate time.

A divergent frequency alone would not settle whether the affected endpoint admits a regular extension, so the decisive calculation uses parallelly propagated transverse curvature. With curvature convention $R(X,Y)Z=\nabla_X\nabla_YZ-\nabla_Y\nabla_XZ-\nabla_{[X,Y]}Z$, the vanishing first transverse metric derivatives on the axis give $R_{AyBy}=-\tfrac12\partial_y^2g_{AB}$ for longitudinal indices $A,B\in\{0,\xi\}$. Contracting with $\ell^A=\omega(1,a)$ and using $w+a=1$ yields $R(\ell,e_y,\ell,e_y)=\beta f'(r)\omega^2/r$, with an identical expression in the other transverse direction.

The unit vectors $e_y=\partial_y$ and $e_z=\partial_z$ are parallelly transported along the axial ray, so these are curvature components in a parallel screen frame, not merely components in a rapidly changing observer frame. Their sum gives $T_{ab}\ell^a\ell^b=2C\beta f'(r)\omega^2/r$, understood as its smooth zero limit in the constant core.

For a nonzero-energy ray approaching a finite-affine endpoint $\lambda_h$, one has $\xi-\xi_h=E(\lambda-\lambda_h)$ and $a=qE(\lambda-\lambda_h)+O((\lambda-\lambda_h)^2)$. Hence $\omega\sim1/[q(\lambda-\lambda_h)]$, and the transverse curvature and null stress behave as $R(\ell,e_y,\ell,e_y)\sim-1/[\kappa h(\lambda-\lambda_h)^2]$ and $T(\ell,\ell)\sim-2C/[\kappa h(\lambda-\lambda_h)^2]$. On either zero-energy generator, the corresponding expressions are exact with $\lambda_h$ replaced by $\lambda_c$. A nondegenerate $C^2$ Lorentzian extension continuing the affine geodesic through such an endpoint would have continuous curvature with finite contractions against its continued tangent and parallel screen vectors, contradicting this divergence.

The result excludes that regular continuation through the affected endpoints; it neither proves inextendibility everywhere nor classifies lower-regularity extensions.

The endpoint obstruction must nevertheless be distinguished from the regular geometry at a zero at finite coordinate time. The longitudinal metric block has determinant $-1$ even when $a=0$, and a different null family can cross those locations. Bounded profile derivatives keep normal-frame curvature and scalar polynomial curvature invariants bounded throughout $M$, but do not bound the parallelly propagated curvature produced by the unbounded affine frequency. An extension based only on the induced axial two-dimensional metric would therefore miss the decisive four-dimensional transverse tidal terms.

The explicit profile also supplies a geometric timescale without introducing a particle-production model. For $\beta=2$, its midpoint gives $h=R+\Delta/2$, $f'(h)=-2/\Delta$, and $\kappa=4/\Delta$. The front generator's affine frequency grows as $\exp(\kappa x^0)$, corresponding to physical e-folding time $\Delta/(4c)$. This is a classical blueshift scale, not a computed quantum emission rate or a proof of quantum instability.

The averaged-energy interpretation now differs sharply from the subluminal result. On each maximal superluminal forward ray, the stress contraction is nonpositive and its integral diverges negatively at every finite-affine endpoint, but none of these rays meets the stipulated two-sided completeness requirement. These are divergent integrals on incomplete rays, not complete-geodesic averaged-null-energy violations. Joining sectors across a zero cannot supply the missing complete ray, because the corresponding asymptotic endpoints fail the regular $C^2$ continuation test.

The backward axial family provides a complete-ray control with the opposite integral sign. Define $b(\xi)=1+\beta-\beta f(|\xi|)=2-a$, so $1\leq b\leq1+\beta$, and write $\ell_-=\omega(\partial_0-b\partial_\xi)$. Its Killing constant is $E_-=\omega b>0$, and $d\xi/d\lambda=-E_-$; the positive upper and lower frequency bounds then give a two-sided infinite affine domain. The axial contraction is $T(\ell_-,\ell_-)=-2C\beta f'(r)\omega^2/r$, yielding $I_-=-4C\beta E_-\int_R^{R+\Delta}f'(r)\big/\{r[1+\beta-\beta f(r)]^2\}\,dr>0$. This integral is finite, and exterior normalization $\omega=1$ fixes $E_-=1+\beta$.

With length-valued affine parameter, these averaged integrals have units $\mathrm{J\,m^{-2}}$; positive affine rescaling changes their magnitude but not their sign or convergence classification. The complete backward result does not establish positivity on all complete geodesics, and the axial classification leaves complete nonaxial integrals undetermined. What changes in the central argument is therefore precise: the specified eternal superluminal target cannot be repaired by treating its affected forward endpoints as removable regular horizons, while its forward negative divergences cannot be used as a complete-ray averaged-energy exclusion.

Any finite-duration replacement must establish its own affine completeness and tidal behavior rather than inherit either a regularity claim or a prohibition from this eternal model.

\section{A wider-cone alternative and the limits of comparison}
A preferred-time, wider-cone field provides a distinct comparison because it changes the propagation structure rather than the geometry carrying an ordinary timelike traveler. In the geometric construction, the carrier remains inside the local metric light cone; the travel advantage concerns a comparison between geometries. The alternative instead retains Minkowski geometry for photons and ordinary matter while assigning a separate scalar field a propagation cone wider than the photon cone. Its superluminality is therefore local relative to the photon sector, not the transport of an ordinary craft along a geometrically shortened route.

The assumptions available for this comparison are a physical global time $t$, positive free-field kinetic energy, a characteristic propagation speed $c_s>c$, and a nonnegative squared-frequency parameter $\omega_0^2\geq 0$. Here $c$ is the photon speed, $c_s$ has units of metres per second, and $\omega_0$ has units of inverse seconds. Future-directed propagation is required to increase the same preferred time. The supplied analysis supports mathematical coherence at the free-field level under these assumptions; the comparison here does not depend on reconstructing an action or reproducing that free-field proof.

The causal role of the preferred time is substantive rather than a coordinate convention. If every permitted propagation segment advances the same global time, a sequence of such segments cannot return to its initial event. This conditional ordering explains how a wider signal cone can be considered without simultaneously granting arbitrary superluminal channels with independently chosen future orientations. It does not establish causal consistency after interactions are added: the common forward-time restriction must remain valid for the complete coupled system, rather than only for the isolated scalar sector.

The principal missing step is an interacting source\textendash{}detector theory. A free propagation sector does not establish that ordinary matter can controllably emit the field, encode information in it, or register its arrival. Those capabilities require specified couplings and a stable coupled evolution; the positive-energy properties of the uncoupled field cannot substitute for that calculation. Conservation must likewise be established for the interacting matter and field system, including their exchanges, and no gravitational completion has been demonstrated.

Even a successful source\textendash{}detector construction would establish signaling rather than payload transport. Moving material requires a further account of payload coupling, its dynamics during transit, and conservation in the resulting interacting theory. No such material-transport sector is supplied. The wider-cone comparison therefore cannot serve as a mechanism for realizing the prescribed geometric corridor, nor can its free-field consistency be counted as evidence for a working travel system.

Empirical compatibility is also unresolved. Preferred-frame effects, additional radiation channels, and radiative corrections are identified as subjects for calculation and observational comparison, not as quantified predictions or established exclusions. No observational dataset or original domain-literature assessment supports numerical preferred-frame limits, a resolution of contested quantum-instability claims, or a claim of literature novelty here. Consequently, the comparison supplies neither empirical acceptance nor empirical rejection of an interacting wider-cone framework.

The alternative changes the scope of the central argument: a requirement on the Einstein source of a prescribed shortcut is not, by itself, a restriction on every possible propagation framework. It replaces the geometric sourcing question with different obligations concerning preferred-time dynamics, interactions, gravity, and material coupling. The resulting distinction is between a conditionally coherent free signal sector and a transport theory\textemdash{}not between two demonstrated ways of achieving superluminal travel.

\section{A conditional corridor and decisive research gates}
A finite-duration transport corridor provides a conditional architecture for testing whether the geometric shortcut can survive source, propagation, quantum, and control requirements together. The proposed architecture separates preparation of the corridor from carriage of information or matter through it, allowing incoming preparation or predeployed infrastructure rather than assuming that a vehicle creates its surroundings on demand. This distinction is consequential: the prescribed transit protocol requires an independently scheduled geometry, not a demonstrated construction triggered by the newly transmitted information.

The research program should therefore treat the corridor as a target for coupled dynamics, not as a device whose feasibility follows from its metric.

The first gate is a specified source action with complete, admissible preparation data. A candidate must reproduce the required energy density, momentum fluxes, and spatial stresses while satisfying the coupled gravitational constraints and matter evolution equations. Defining a tensor from the Einstein tensor, even when it has the necessary geometric conservation property, does not supply those dynamics or establish a preparable state. Advancement through this gate requires a well-posed source\textendash{}gravity initial-value problem rather than another identification of required stress with realizable matter.

The source exclusions already obtained determine which candidates should not proceed unchanged. Under the stated assumptions, neither the canonical scalar nor the regular minimally coupled single-scalar $P(X,\phi)$ theory with positive high-frequency physical energy and a common, nondegenerate matter\textendash{}gravity Cauchy slicing can support this exact target. Changing preparation alone cannot repair that local incompatibility; a proposal must change the action, the target geometry, or the regular-propagating-mode assumptions and then justify the replacement.

Independently, an autonomous theory with a unique Cauchy development cannot produce the pulse from complete vacuum data, so incoming fields, nontrivial initial source data, or another explicit departure from those hypotheses must be identified.

The second gate is an explicitly finite-duration propagation calculation, addressing whether switching changes the obstruction found in the eternal geometry. The proposed discriminator is to test whether a finite time integral of the relevant spatial velocity gradient bounds affine frequency through the pulse. This remains an unproved criterion, not an established regularity theorem or a demonstrated repair of the eternal endpoints. The calculation must determine affine evolution and parallelly propagated curvature for the finite episode rather than importing the stationary conserved quantities or treating the eternal pathology as a removable coordinate horizon.

A favorable frequency bound would answer only part of the propagation question. Complete-ray energy integrals require separate checks of affine completeness, normalization, and convergence, including nonaxial rays before any global averaged-null-energy conclusion is drawn. Tidal exposure likewise requires curvature evaluated along the relevant trajectories; bounded frequency alone would not establish acceptable exposure for a transported system. The proposed failure criteria at this gate are an unresolved finite-affine endpoint obstruction, uncontrolled tidal curvature, or failure of an energy constraint shown to apply to the specified model\textemdash{}not merely a negative integral over an incomplete ray.

The third gate is a model- and state-specific quantum calculation on the finite-duration configuration. It requires specified fields, a quantum state, renormalization conditions, and boundary conditions before quantum energy inequalities, renormalized stress, and stress fluctuations can be evaluated. Those results must then enter a self-consistent solution or a controlled bound on backreaction. Classical negative null stress and classical blueshift do not, by themselves, establish either a universal quantum exclusion or quantum safety.

For this gate, the decisive question is whether the proposed state and its backreaction remain compatible with the source and geometry required for the transit. A calculation on a fixed background would be an intermediate test, not a substitute for examining how the calculated stress changes that background. Advancement would require a justified regime in which the mean stress and its fluctuations support a controlled description; failure to find such a regime would reject that quantum realization. Neither a specified state nor this backreaction calculation is presently supplied.

The fourth gate is causal preparation and operational control through formation, transit, switching, and shutdown. A proposed realization must determine which source regions can receive which control signals, particularly at the leading wall, rather than equating a prescribed time-dependent profile with an executable control law. Predeployment is an architectural alternative to onboard creation, but it still requires admissible preparation data and a demonstration that the intended sequence follows from them. The resulting evolution must also retain an exterior-clock arrival advantage once formation and shutdown are included; the existing kinematic timing comparison does not establish that dynamical sequence.

Control must be assessed together with perturbative and nonlinear stability of the coupled system. The local source screen is a necessary filter, not evidence that a replacement theory will preserve the corridor under finite disturbances or interactions. Similarly, the temporal function available for one prescribed arrangement does not settle the causal structure of several corridors, relative motions, or repeated trips. Each proposed assembly must therefore pass a new causal analysis; a loop-free isolated configuration cannot serve as blanket authorization for a transport network.

The final gate is confrontation with observations for a sufficiently specified theory. No observational dataset or original domain-literature assessment is supplied here, so empirical compatibility, quantitative preferred-frame limits, contested quantum-instability claims, and literature novelty remain unassessed. For the wider-cone alternative, reaching this gate also requires stable interacting source, detector, and gravitational sectors, with payload coupling, conservation, and radiative control; free-field signaling is not material transport.

Observational tests must consequently be tied to the interactions and propagation properties of an explicit candidate rather than to the mere availability of a wider cone or a prescribed shortcut metric.

These gates change the next useful target from geometric optimization to two discriminating tasks: finding a source that escapes the established exclusions without sacrificing a well-posed physical evolution, and determining whether a finite episode avoids the eternal propagation obstruction. Either task can reject a particular source\textendash{}target combination without establishing a prohibition on every effective-superluminal framework. Only their joint success would justify advancing the conditional corridor to quantum and operational transport tests; the geometric shortcut alone no longer carries that inference.

\section{Limitations}
\begin{itemize}
\item The source selection is bounded and omits substantial underlying notes. It supports a scoped synthesis, not an exhaustive review or independent reconstruction of every proof.
\item The conditional geometry and scalar exclusions concern specified metrics, source classes, and classical assumptions. They do not exclude every possible spacetime shortcut or replacement theory.
\item The superluminal endpoint obstruction concerns the specified eternal geometry and forward axial family. It neither proves a finite-duration obstruction nor settles averaged energy on all complete geodesics.
\item The spherical energy estimate is a chosen-slice integral, not an asymptotic mass or an engineering energy budget. The supplied preview does not justify a numerical device specification.
\item The alternative framework is available only as a reported free-field comparison; its omitted equations prevent an equally detailed mathematical assessment.
\item No measured observation, completed simulation, specified quantum-state backreaction calculation, or inspected domain-literature review establishes physical viability or observational compatibility.
\end{itemize}
\section{Open obligations}
\begin{itemize}
\item Supply a matter or replacement-gravity action and solve its coupled constraints and evolution for a source matching the required density, pressures, and momentum fluxes. Establish material conservation and a common well-posed evolution rather than relying on geometric conservation alone.
\item Specify complete nonvacuum initial data, incoming preparation, or boundary infrastructure. Demonstrate formation without invoking spontaneous evolution from complete Minkowski vacuum data.
\item Complete the finite-duration affine calculation proposed in C19, including all relevant ray families and switching regions. Determine whether the eternal blueshift pathology persists or is avoided under explicit bounds.
\item Determine averaged null-energy integrals on complete nonaxial geodesics and applicable finite-duration rays; do not use incomplete forward-ray integrals as substitutes.
\item Choose quantum fields and states, calculate renormalized stress and fluctuations, identify applicable quantum energy bounds, and solve or bound self-consistent backreaction.
\item Establish coupled perturbative and nonlinear stability, tidal exposure, leading-wall control, shutdown, and safe arrival. Distinguish externally prepared corridors from a self-contained craft.
\item Analyze interacting and repeated configurations for causal consistency; the single-metric temporal function is insufficient for arbitrary networks.
\item For the preferred-time alternative, provide the omitted free-field equations and construct stable source, detector, matter, and gravity interactions with explicit conservation laws and payload coupling.
\item Inspect actual observational records and original physical literature before assigning preferred-frame limits, empirical compatibility, contested quantum-instability conclusions, or novelty.
\item Define quantitative acceptance and rejection thresholds only after a concrete source theory, quantum state, operational protocol, and observational comparison have been specified.
\end{itemize}
\section{Conclusions}
The prescribed finite-duration geometry supplies a conditional counterexample to the assertion that an exterior-clock shortcut must involve locally superluminal carrier motion or a closed causal curve. It does not supply the physical mechanism required to realize that geometry.

The source analysis is restrictive rather than constructive: the target requires strict null-energy violation, while the canonical scalar and the stated regular positive-energy single-scalar extension cannot provide it. The eternal superluminal model also has a genuine axial endpoint pathology that cannot be repaired by treating the affected endpoints as removable coordinate horizons.

Neither a universal prohibition nor a working effective-superluminal transport framework is established. The supported outcome is a conditional corridor concept and a staged research program, with finite-duration propagation, replacement-source dynamics, quantum backreaction, control, and empirical compatibility serving as decisive gates.

\section*{Claim and source notes}
These notes identify the saved evidence for claims used in the article. The complete claim register and reasoning are retained in the paper record.
\begin{sloppypar}
\small
\begin{description}
\item[C11] For the specified eternal constant-speed metric with $\beta>1$ and the explicit smooth monotone spherical profile, every maximal future-directed forward axial null geodesic is affinely incomplete in at least one direction. Its finite-affine endpoint has divergent parallelly propagated tidal curvature and cannot be continued through that endpoint by a regular $C^2$ extension. \textit{Status: }established.
\textit{Evidence: }[Research session 1, artifact 6]
\item[C12] For the same eternal superluminal target, the backward axial null family is complete and has a finite positive averaged null-energy integral. The incomplete forward family does not furnish a complete-geodesic averaged-null-energy test. \textit{Status: }established.
\textit{Evidence: }[Research session 1, artifact 6]
\item[C13] A preferred-time wider-cone free-field model provides a formally different comparison with geometric shortcuts, but no interacting source\textendash{}detector, gravity, or material-transport sector has been established for it. \textit{Status: }limitation.
\textit{Evidence: }[Research session 1, artifact 1], [Research session 1, artifact 2]
\item[C14] The prescribed geometry and its required Einstein tensor do not establish a preparable material source, stable coupled evolution, or operational transport mechanism. \textit{Status: }limitation.
\textit{Evidence: }[Research session 1, artifact 1], [Research session 1, artifact 5], [Research session 1, artifact 6]
\item[C15] Quantum energy constraints, renormalized stress, fluctuations, and self-consistent backreaction remain undetermined for a specified field and state on a finite-duration shortcut. \textit{Status: }unresolved.
\textit{Evidence: }[Research session 1, artifact 1], [Research session 1, artifact 2], [Research session 1, artifact 6]
\item[C16] Admissible formation, switching, leading-wall control, shutdown, tidal safety, nonlinear stability, and causal consistency of combined transport configurations remain unresolved. \textit{Status: }unresolved.
\textit{Evidence: }[Research session 1, artifact 1], [Research session 1, artifact 4], [Research session 1, artifact 5]
\item[C17] Compatibility with existing observations, quantitative preferred-frame limits, contested quantum-instability literature, and literature novelty are not assessed by the supplied evidence. \textit{Status: }limitation.
\textit{Evidence: }[Research session 1], [Research session 1, artifact 4], [Research session 1, artifact 5]
\item[C18] A finite-duration, possibly predeployed transport corridor is a conditional conceptual design whose advancement requires successive source, propagation, quantum, control, and observational gates; it is not a demonstrated drive. \textit{Status: }conjecture.
\textit{Evidence: }[Research session 1, artifact 1], [Research session 1, artifact 4], [Research session 1, artifact 5], [Research session 1, artifact 6]
\item[C19] An explicitly finite-duration propagation calculation is proposed to test whether a finite integrated spatial velocity gradient bounds affine frequency and avoids the eternal geometry's endpoint pathology. \textit{Status: }conjecture.
\textit{Evidence: }[Research session 1, artifact 6]
\item[C10] For $0<\beta<1$ and a smooth monotone spherical profile, the complete forward axial null family has a strictly negative averaged null-energy integral under the stated exterior normalization. \textit{Status: }established.
\textit{Evidence: }[Research session 1, artifact 5]
\end{description}
\end{sloppypar}
\section*{Paper writing source material used}
These sources informed manuscript organization and writing methods, not the scientific evidence or findings.
\begin{itemize}
\item IOP \textemdash{} Writing Your Journal Article (discipline): \url{https://publishingsupport.iopscience.iop.org/publishing-support/authors/authoring-for-journals/writing-journal-article/}
\item PLOS ONE \textemdash{} Manuscript and Reporting Guidelines (general): \url{https://journals.plos.org/plosone/s/submission-guidelines}
\item EASE \textemdash{} Guidelines for Authors of Scientific Articles (full text) (general): \url{https://pmc.ncbi.nlm.nih.gov/articles/PMC4130677/}
\end{itemize}
\section*{Authorship note}
AI-assisted synthesis from saved research; not independently peer reviewed.
\end{document}
